\thispagestyle{plain}
\section*{译本体例}
\addcontentsline{toc}{section}{译本体例}

\subsection*{一、编号}
原书每一讲独立编号，讲内各节标为“1.”“2.”……；每节之内，一切带编号的对象（口号、定义、定理、命题、例、练习、构造、注，以及以“2.1.”开头的普通段落）共用同一个计数器，因此同一节内编号连续。本译本完整保留这一体例：节号与对象号均与原书一致，正文中的交叉引用（如“由定理 2.2”）也按原书编号给出。

\subsection*{二、译注}
译注以左侧竖线的小字排于正文之后，随文就近安置。译注分三类：
\begin{itemize}[leftmargin=2.2em,itemsep=0.15em]
\item 术语说明：解释译名取舍，或在首见处补出通行英文；
\item 背景补充：补足理解所需、而原讲义从略的定义或事实；
\item 疑点讨论：原文表述含混、笔误或前后不一致之处。
\end{itemize}
每条解释性译注末尾标注依据等级：\textbf{A}（可核对的公开文献）、\textbf{B}（同类教材通例）、\textbf{I}（据上下文推断），或“译者按（无外部依据）”。

\subsection*{三、术语}
术语首见时以“中文（English）”给出，其后一律用中文。全书术语以《工程文件/work/glossary.tsv》为准；几个关键取舍如下：
\begin{center}
\small
\begin{tabular}{lll}
\toprule
原文 & 译名 & 说明\\
\midrule
$\infty$-category & $\infty$-范畴 & 默认指 $(\infty,1)$-范畴，与原书一致\\
quasi-category & 拟范畴 & 不译“准范畴”\\
simplicial set & 单纯集 & \\
Kan complex & Kan 复形 & 人名保留\\
model category & 模型范畴 & \\
fibration / cofibration & 纤维化 / 上纤维化 & \\
homotopy coherent & 同伦连贯 & 名词作“同伦连贯性”\\
$\infty$-operad & $\infty$-操作子 & 不译“运算子”\\
straightening / unstraightening & 平展 / 去平展 & Lurie 的 straightening 函子\\
$t$-structure & $t$-结构 & \\
presentable & 可表现 & \\
\bottomrule
\end{tabular}
\end{center}

\subsection*{四、符号}
\begin{center}
\small
\begin{tabular}{llll}
\toprule
记号 & 读法 & 记号 & 读法\\
\midrule
$\cC$, $\cD$ & 范畴 & $\mathrm{Set}$ & 集合范畴\\
$\sSet$ & 单纯集范畴 & $\Cat$ & 小范畴范畴\\
$\Spc$ & 空间（$\infty$-群胚） & $\Sp$ & 谱\\
$\Delta^n$ & $n$ 单形 & $\partial\Delta^n$ & 边界\\
$\Lambda^n_i$ & 角（horn） & $\Hom_{\cC}(x,y)$ & 态射空间\\
$\Fun(\cC,\cD)$ & 函子范畴 & $\Map_{\cC}(x,y)$ & 映射空间\\
$\colim$, $\lim$ & 余极限、极限 & $\hocolim$, $\holim$ & 同伦余极限、同伦极限\\
$\cC_{/x}$, $\cC_{x/}$ & 切片、余切片 & $\cC^{\ot n}$ & $n$ 重张量积\\
$\ho\cC$ & 同伦范畴 & $\cC_{\sim}$ & 局部化\\
\bottomrule
\end{tabular}
\end{center}

\subsection*{五、图}
原书共五幅图（Figure 1–5），均为单形与角的示意图。本译本按原意用 Ti\emph{k}Z 重绘，图号与原书一致。

\subsection*{六、引用}
原书随讲给出参考文献，体例为作者缩写加年份（如 [Lur09]）。本译本保留原标签，不作统一书目合并；各讲末尾的“参考文献”即原书该讲所列条目。
